\documentclass[border=4pt]{standalone}
\usepackage[siunitx]{circuitikz}

\begin{document}
\begin{circuitikz}[font=\small,line width=0.9pt]
  \ctikzset{tripoles/thickness=1,logic ports/thickness=1,
    flipflops/thickness=1}
  \tikzset{
    block/.style={draw, minimum width=1.8cm, minimum height=1.25cm,
      align=center, font=\small}
  }

  \node[font=\bfseries] at (4.0,5.4) {(a) Functional D-latch feedback model};
  \node[block] (STORE) at (6.4,3.4) {inverting-pair\\storage delay};
  \coordinate (DPORT) at (2.5,3.75);
  \coordinate (FBPORT) at (2.5,3.05);
  \coordinate (MUXOUT) at (4.3,3.4);
  % At x=3.4, the select lead meets the MUX's lower sloped edge at y=2.85.
  \coordinate (MUXSEL) at (3.4,2.85);
  \coordinate (DIN) at ($(DPORT)+(-1.3,0)$);
  \coordinate (ENIN) at ($(MUXSEL)+(0,-1.0)$);
  \coordinate (QOUT) at ($(STORE.east)+(0.9,0)$);
  \coordinate (QJ) at ($(STORE.east)+(0.6,0)$);
  \coordinate (FBY) at (3.4,1.42);
  \coordinate (FBTURN) at ($(FBPORT)+(-0.55,0)$);
  \draw (2.5,2.70) -- (4.3,3.00) -- (4.3,3.80)
    -- (2.5,4.10) -- cycle;
  \node[font=\small] at (3.35,3.4) {MUX};
  \node[font=\scriptsize] at (2.72,3.75) {1};
  \node[font=\scriptsize] at (2.72,3.05) {0};
  \draw (DIN) node[ocirc, label={[font=\small]left:{$D$}}] {}
    -- (DPORT);
  \draw (ENIN) node[ocirc, label={[font=\small]below:{$EN$}}] {}
    -- (MUXSEL);
  \draw (MUXOUT) -- (STORE.west);
  \draw (STORE.east) -- (QOUT)
    node[ocirc, label={[font=\small]right:{$Q$}}] {};
  \draw (QJ) -- (QJ |- FBY) -- (FBTURN |- FBY)
    -- (FBTURN) -- (FBPORT);
  \fill (QJ) circle (1.4pt);
  \node[font=\small, below, align=center] at (4.9,0.72)
    {$EN=1$: input path selected\quad $EN=0$: feedback selected};

  \draw[densely dashed] (9.1,0.7) -- (9.1,5.7);

  \node[font=\bfseries] at (14.0,5.4)
    {(b) positive-edge D flip-flop};
  \node[latch,flipflop def={n3=1},scale=0.95] (L1) at (12.1,3.4) {};
  \node[latch,scale=0.95] (L2) at (15.7,3.4) {};
  \node[above] at (L1.north) {L1: low-phase latch};
  \node[above] at (L2.north) {L2: high-phase latch};
  \coordinate (D2) at ($(L1.pin 1)+(-1.2,0)$);
  \coordinate (Q2) at ($(L2.pin 6)+(1.2,0)$);
  \coordinate (C1) at ($(L1.pin 3)+(-0.35,0)$);
  \coordinate (C2) at ($(L2.pin 3)+(-0.35,0)$);
  \coordinate (CLKIN) at (10.1,1.25);
  \coordinate (CLKEND) at ($(C2 |- CLKIN)+(0.75,0)$);
  \draw (D2) node[ocirc, label={[font=\small]left:{$D$}}] {} -- (L1.pin 1);
  \draw (L1.pin 6) -- (L2.pin 1);
  \draw (L2.pin 6) -- (Q2)
    node[ocirc, label={[font=\small]right:{$Q$}}] {};
  \draw (CLKIN) node[ocirc, label={[font=\small]below:{$CLK$}}] {}
    -- (CLKEND);
  \draw (C1 |- CLKIN) -- (C1) -- (L1.pin 3);
  \draw (C2 |- CLKIN) -- (C2) -- (L2.pin 3);
  \fill (C1 |- CLKIN) circle (1.4pt);
  \fill (C2 |- CLKIN) circle (1.4pt);
  \node[font=\small, below, align=center] at (14.0,0.45)
    {The bubble makes L1 active low; L2 is active high.\\
     Only one latch is transparent in either settled clock phase.};
\end{circuitikz}
\end{document}
